% =====================================================================
% paper/sections/related_work.tex
% §2 Related Work（rev.1, 2026-05-01）
%
% 写作纪律（outline §2）：
%   - 3 段：(2.1) offline RL with behavior anchoring  →  锁定 method 谱系
%           (2.2) underwater sim2real navigation       →  锁定 application 谱系
%           (2.3) asymmetric / privileged actor-critic →  锁定 finding 2 的方法学谱系
%   - 每段 4-6 句，引用 3-5 篇
%   - 关键是讲清三条线各自的局限，让审稿人接受为何要写这篇 paper
%   - 不与 §1 / §6 重复机制论证；只交代 method-vs-prior-work 的边界
%
% 锚点：sec:related_work（被 intro.tex 1 处 forward-ref）
% =====================================================================

\section{Related Work}\label{sec:related_work}

\paragraph{Offline RL with Behavior Anchoring.}
本文工作落在 ``policy-constraint'' offline RL 谱系上。该路线的早期工作 BCQ \citep{fujimoto2019bcq} 通过 perturbed VAE 限制 action 在 dataset support 内；后续 CQL \citep{kumar2020cql} 在 critic 端施加 OOD action 的 Q 值惩罚；IQL \citep{kostrikov2022iql} 通过 expectile regression 把 critic 学习与 action 选择解耦从而完全避免 OOD action 评估；AWAC \citep{nair2020awac} 用 advantage-weighted BC 联系 offline / online 阶段。在这一谱系中，TD3+BC \citep{fujimoto2021td3bc} 与 ReBRAC \citep{tarasov2023rebrac} 因其 ``minimal recipe'' 风格自成一支 —— 前者把约束简化为 actor loss 中的单侧 BC penalty $\alpha \cdot \mathbb{E}[\|\pi_\theta(s) - a\|^2]$，后者在此之上扩展为 actor / critic 双侧 BC penalty 与 critic LayerNorm。我们的 ReBRAC-Q 复用 ReBRAC 的 dual penalty + LN minimal recipe，但在 actor loss 上保留 TD3+BC 的 $Q$ 归一化（详见 \S\ref{subsec:method_diff}）—— 这一组合在文献中尚未被系统消融，本文 finding 3 与 finding 4 因此是对该 minimal recipe 内部各 component 独立必要性的首份机制级拆解。

\paragraph{Underwater Sim2real Navigation.}
水下 AUV 的强化学习导航研究多落在 fluid-aware 控制与 wake-aware 路径规划两个子方向 \citep{underwater_rl_placeholder} —— 前者在均匀或低维流场假设下学习反馈控制，后者用 model-based / planning 在已知流场地图上构造路径。本文与该方向的核心区别有三：(i) 我们使用 D2Q9 TRT-LBM 离线生成的真实 Kármán 涡街尾迹场（\S\ref{subsec:setup_task}），而非解析叠加或低维近似；(ii) 我们的 deployable sensor 协议（\texttt{s0} = 标配 1200 kHz DVL 水追踪单点 \citep{fossen2011handbook}）严格等同真实 REMUS-100 部署侧能获得的最弱信号，没有任何上游 / 侧向 / 整流壳积分级别的 simulator-only 通道；(iii) 我们关心的是离线策略训练而非在线探索 —— 训练阶段一切在 simulator 中发生，部署阶段策略只读 deployable sensor。在这三条边界共同作用下，本文是据我们所知第一份在 Kármán 尾迹场 + 真实 DVL 协议下系统评估 offline RL 的工作。

\paragraph{Asymmetric / Privileged Actor-Critic.}
非对称 actor-critic 框架 \citep{pinto2017asymmetric} 是机器人 sim2real 的 standard recipe：训练时给 critic 注入 privileged 信息（如完整状态、ground-truth 物体姿态、特权感知通道）以加速 value 学习，测试时仅 actor 被部署、不依赖任何 privileged 信号。该 recipe 的隐含假设是 ``critic 越聪明，actor 越聪明''。本文 finding 2（\S\ref{subsubsec:finding_2}）对该假设给出一个边界条件：在我们的 ReBRAC-Q + s0 sensor + Kármán 尾迹场 setup 下，critic 端的 privileged 通道（hull-integral $[u_{\mathrm{eq}}, v_{\mathrm{eq}}]$）\emph{不再}贡献 mean uplift，仅救援 outlier seed —— actor-side BC penalty 已把 mean 抬到 privileged-critic 协议的同一上界（Welch's $p = 0.92$）。这一逆向结果不否定 asymmetric AC 的一般有效性，而是给出一个具体的 ``actor 端约束已饱和'' 制度：在该制度下，加 privileged 通道只能提供冗余的 critic 信息，无法继续抬 mean。
